\section{两座桥的条件账本与结论}\label{sec:bridges}
\subsection{解析桥：结构已落实的部分与未支付的估计}
图像 Bridge A 从岛②通往岛⑥。仓库的 Gate A 与 Gate B 都是这座桥内部的解析义务，并非地图上的两座不同红桥。
表~\ref{tab:ledger}将本稿使用的数量放在共同尺度下。
\begin{table}[htbp]
\centering\small
\begin{tabularx}{\linewidth}{@{}lXl@{}}
\toprule
环节 & 共同尺度上的内容 & 状态\\\midrule
实际算子 & $Q^2/\log Q\ll\norm A\ll Q^2$；全部实际行能量 $\asymp HQ^2/\log Q$ & 已证限定结论\\
实际系数 & 范数乘积法不能单独支付固定幂节省 & 已证方法障碍\\
边界重组 & $xQ^2x^{-11/64+o(1)}$ & 已支付此部分\\
临界数组输入 & 合法双侧短数组因子 $q^{31/32+o(1)}$ & 外部定理\\
双侧附接 & 条件~\ref{cond:compression}及全部残差 & 未证\\
全标量目标 & $|\Cx|\ll xQ^2x^{-\delta+o(1)}$, $\delta>1/400$ & 未证\\
解析 Gate A & 原完整壳层有符号折叠与正余量 $\eta_A$ & 另行未证\\
\bottomrule
\end{tabularx}
\caption{指数账本。算子上下界仍有对数差距。结构恒等式、方法障碍与目标上界不是同一种证据。}\label{tab:ledger}
\end{table}
\paragraph{导入的条件接口与参照指数。}
为明确这个账本估计什么，记
\[
S_x^{\rm phys}=\sum_{t\in I_x}\beta(t)w(t),\qquad
K_* =\sum_{q\in\QQ}q\frac{q-2}{q-1}=x^{2/3+o(1)}.
\]
$S_x^{\rm phys}$ 不是第~\ref{sec:arithmetic}节的 $S_2(x)$。
V58 的源折叠标量 $A_*$ 与误差标量 $E_*$ 给出以下导入接口\cite{archivev58}：
\[
S_x^{\rm phys}=\frac{A_*-\Cx+E_*}{K_*},\qquad
|E_*|\ll x^{143/96+o(1)}.
\]
这里 $A_*$ 是复标量，不是矩阵 $A$；其来源定义为全壳层折叠 $\sum_q qP_q$，
$E_*=\sum_q qE_q$，逐行 $P_q,E_q$ 的构造和误差证明由所引源接口提供。
本稿将这组恒等式及误差界作为导入前提，不重新构造其上游行数据。
解析 Gate A 的尚未证明假设则明确为
\[
|A_*|\ll x^{1997/1200-\eta_A+o(1)},\qquad \eta_A>0.
\]
在该假设及全标量目标 $|\Cx|\ll x^{5/3-\delta+o(1)}$ 下，除以 $K_*$ 后三个输出指数分别是
\[
\frac{399}{400}-\eta_A,\qquad 1-\delta,\qquad \frac{79}{96}.
\]
因此得到条件性输出 $|S_x^{\rm phys}|\ll x^{399/400-\eta+o(1)}$，其中
\begin{equation}
0<\eta<\min\{\eta_A,\delta-1/400,419/2400\}.\label{eq:terminal}
\end{equation}
式~\eqref{eq:terminal}正是三项到参照指数 $399/400$ 的余量：
$(399/400)-(1-\delta)=\delta-1/400$，
$(399/400)-(79/96)=419/2400$。
这是源契约的指数推导\cite{archivev58,archivev59,archivephysical}，不是本稿重新独立证明了全部上游到 TPC 的蕴含链。
尤其，单独证明本稿标量的节省也不能省略 Gate A 或历史上游结构附接。
临界数组路线若成功，$\delta=1/96-\ell$，所以严格需要 $\ell<19/2400$；
若采用其他内部估计，边界重组将本稿能认证的全节省限制到 $\min(\delta_{\rm int},11/64)$。

\subsection{动力学桥：度量定理与指定点之间}
图像 Bridge B 从岛⑤通往岛⑥。它已经有精确筛事件编码、质量估计和
定理~\ref{thm:haar}的度量命中定理；缺的是将结论落实到指定种子 $0$。
命题~\ref{prop:seed}将一种足够的缺口写成柱内振荡与点值代价，
但并没有给这两项所需的联合上界。岛⑦若提供表示，也必须回到这一点值要求或另一个同样精确的算术读出条件。

\subsection{统一稿得到什么}
七岛的价值在于区分研究对象和转移任务，而不是把历史编号排序成一条已闭合的证明链。
基础算术给出局部一致性；固定原子与出现项层锁定系数、支持和前缀；
解析层给出实际算子结构、增长维数障碍和可支付边界；动力学层给出实际移动靶的度量命中与精确筛几何。
这些结论具有各自明确的数学内容，既不应混成 TPC 证明，也不应一概归为玩具模型。

本稿是一篇以现有结果为对象的研究综合，尚非完成查新与投稿定位的最终论文。
下一版若面向外部读者，应进一步压缩历史接口、补充更系统的外部相关工作，并明确作者与目标期刊。
但数学上的关键变化不能靠文字整理产生：解析桥仍需实际系数敏感的有符号估计，动力学桥仍需指定算术种子的回接。本文未证明孪生素数猜想。
